% MT010 / explanation — canonical editable content.
先看图的结构：绿色负类在内层，红色正类分布更宽、数量更多，D靠近内层区域。题目让两种模型判断同一个D，是要比较它们决定边界的依据。


\textbf{Gaussian Bayes比较概率证据。} 为两类分别拟合高斯分布，再乘类别先验。某类样本多，估出的先验通常更大；但某类分布很宽，同样的位置上密度又可能更低。因此不能只看D离哪个中心近，也不能只看哪类点多。原解在其示意拟合下预期D属于红色\(+1\)类。原图两类均为圆点，本题建议保留彩色。


\textbf{二次核SVM寻找特征空间的间隔分界。} 二次特征使原空间可以出现曲线，例如把内层点围住、把外圈留在另一侧。配套草图把D放在绿色内层一侧，因此预期为\(-1\)。它没有先拟合两个高斯密度。


精确的高斯分类需要各类均值、协方差与先验；精确的SVM曲线还依赖\(C\)与实际点坐标。原图没有给齐这些量，因此两类结果是配套解答的示意判断。


选读：Gaussian Bayes的矩阵分数。把待判断点D的特征写成向量，\(\mu_c\)是\(c\)类各维均值，\(\pi_c\)是类别先验。协方差矩阵\(\Sigma_c\)的对角线记录各维方差，其他位置描述两维怎样一起变化。假定它正定且可逆，两类可比较



\[
g_c(D)=\log\pi_c-\frac12\log|\Sigma_c|
-\frac12(D-\mu_c)^T\Sigma_c^{-1}(D-\mu_c),
\]


其中省略了两类共同的常数。\(\Sigma_c^{-1}\)是与\(\Sigma_c\)相乘得到单位矩阵的逆矩阵；\(|\Sigma_c|\)是行列式，这里反映分布在各方向上的总体宽窄。最后的二次型按各方向的波动和关联调整距离。


若各维独立，\(\Sigma_c\)只有对角线上的方差\(v_j\)，最后的距离项便是\(\sum_j(D_j-\mu_{cj})^2/v_j\)，而\(\log|\Sigma_c|=\sum_j\log v_j\)。这就回到\StudyRef{Answers}{mt046}{MT046}的逐维高斯分数；完整矩阵写法还允许各维相关。


\textbf{出处：} 2020B期中Quiz Q10，原卷第5页（10分）


\textbf{答案依据：} 2020B期中 Q10 随卷／配套参考答案；条件与解释经复核。
